\documentclass[10pt]{article}
\usepackage[margin=0.82in]{geometry}
\usepackage[T1]{fontenc}
\usepackage{lmodern}
\usepackage{microtype}
\usepackage{amsmath,amssymb,amsthm}
\usepackage{xcolor}
\usepackage[hidelinks]{hyperref}

\definecolor{statusblue}{RGB}{32,76,122}
\newtheorem{theorem}{Theorem}
\newcommand{\Czero}{C_0(\mathbb R)}
\newcommand{\Bdot}{\dot{\mathcal B}'(\mathbb R)}
\newcommand{\indlim}{\mathop{\mathrm{ind}}\limits}

\title{The Universal Extrapolation Space of the Shift on $C_0(\mathbb R)$}
\author{Literature-implied full resolution of arXiv:1604.00763}
\date{9 August 2026}

\begin{document}
\maketitle

\noindent
\colorbox{statusblue!10}{\parbox{0.95\linewidth}{
\textbf{Status.} \texttt{literature\_implied\_answer}.  The conjecture is
fully resolved.  The decisive topological representation is published in a
later paper; the short cofinality identification with the extrapolation scale
is made explicit here.
}}

\section*{Source question}

Bargetz and Wegner ask in Section 4, PDF page 8 of
\cite{BW}, whether the universal extrapolation space for the shift semigroup
on $X=\Czero$ coincides with the space $\Bdot$ of distributions vanishing at
infinity.  In their notation, the universal space is the completion of the
locally convex inductive limit of the negative extrapolation spaces.

As usual in the extrapolation construction, shift the generator by a
resolvent point.  Write $D=d/dx$ in the distributional sense and set
\[
 A=I-D.
\]
Both $I-D$ and $I+D$ have bounded inverses on $\Czero$ (indeed, the inverses
are convolution/translation averages against one-sided exponential kernels).
The shift does not change the domains $D(A^n)=C_0^n(\mathbb R)$ or the
associated universal scale up to its canonical equivalent realization.

\section*{Decisive supporting theorem}

Bargetz, Nigsch, and Ortner prove on PDF page 6 of \cite{BNO} the
\emph{topological} identity
\begin{equation}\label{eq:bno}
 \Bdot
 =\indlim_{m\to\infty}(I-D^2)^m\Czero,
\end{equation}
where every Banach step carries the final norm transported from $\Czero$ by
$(I-D^2)^m$.  They also explain why this is a Hausdorff inductive limit and
why the displayed equality is topological, not merely algebraic.

\section*{Identification}

\begin{theorem}
Let $(T(t))$ be the translation semigroup on $\Czero$, and form its standard
extrapolation scale using the invertible shifted generator $A=I-D$.  Then
\[
 \indlim_{n\to\infty}X_{-n}\cong\Bdot
\]
topologically.  In particular, the inductive limit is Hausdorff and complete,
and the universal extrapolation space satisfies
\[
 X_{-\infty}\cong\Bdot.
\]
\end{theorem}

\begin{proof}
First realize each abstract extrapolation step inside the distributions.  By
definition, $X_{-n}$ is the completion of $X=\Czero$ for
\[
 \|x\|_{-n}=\|A^{-n}x\|_\infty.
\]
If $(x_k)$ is Cauchy for this norm, then $A^{-n}x_k$ converges in $\Czero$ to
some $f$, while $x_k=A^n(A^{-n}x_k)$ converges distributionally to $A^nf$.
Conversely, density of $D(A^n)$ in $\Czero$ gives every $A^nf$ this way.
Thus, isometrically up to the harmless chosen graph-norm convention,
\begin{equation}\label{eq:E}
 X_{-n}=E_n:=A^n\Czero,
 \qquad \|A^nf\|_{E_n}=\|f\|_\infty.
\end{equation}

Put $C=I+D$ and $F_m=(I-D^2)^m\Czero=A^mC^m\Czero$, with its transported
Banach norm.  We show that $(E_n)$ and $(F_m)$ are cofinal by continuous
inclusions.

For $A^nf\in E_n$, bounded invertibility of $C$ gives
\[
 A^nf=A^nC^n(C^{-n}f)\in F_n.
\]
Hence $E_n\hookrightarrow F_n$ continuously.  In the other direction, note
that $C=2I-A$, so
\[
 R:=A^{-1}C=2A^{-1}-I
\]
is bounded on $\Czero$.  For $A^mC^mf\in F_m$, commutativity of these
constant-coefficient operators yields
\[
 A^mC^mf=A^{2m}R^mf\in E_{2m}.
\]
Thus $F_m\hookrightarrow E_{2m}$ continuously.

These two families of inclusions induce mutually inverse continuous identity
maps between the two locally convex inductive limits.  Consequently
\[
 \indlim_n X_{-n}
 =\indlim_n E_n
 \cong\indlim_m F_m
 \overset{\eqref{eq:bno}}{=}\Bdot.
\]
The right-hand side is complete in the topology used in \cite{BNO}.  Therefore
the completion appearing in the definition of $X_{-\infty}$ does not enlarge
the limit, which proves the final assertion.
\end{proof}

\section*{Literature status and scope}

This is a full answer to the concrete $C_0(\mathbb R)$ shift conjecture, not a
solution of the source paper's broader question for arbitrary non-reflexive
Banach spaces.  The supporting authors cite \cite{BW}, but \cite{BNO} does not
state that representation \eqref{eq:bno} resolves the conjecture.  The
relation is therefore an identified implication of published literature.

The bounded status check searched the run's registry, solution, attempt, and
proof-gap indexes for the arXiv id, title, and core terms; exact-title and
exact-phrase web searches; and the OpenAlex citation list for the published
source article (DOI \texttt{10.1016/j.jmaa.2017.11.042}).  OpenAlex listed one
citing work, namely \cite{BNO}.  No paper explicitly announcing a resolution
was found.  Novelty is not claimed because \eqref{eq:bno} is already a
published topological theorem.

\section*{Human-review recommendation}

Verify the standard realization \eqref{eq:E} and the convention that the
shift generator is replaced by the resolvent shift $I-D$.  The decisive
finite-step checks are the two inclusions
$E_n\hookrightarrow F_n$ and $F_m\hookrightarrow E_{2m}$; both are explicit
above.  If those conventions match the source's universal construction, the
resolution is complete.

\begin{thebibliography}{9}
\bibitem{BW}
C. Bargetz and S.-A. Wegner,
\emph{Pivot duality of universal interpolation and extrapolation spaces},
J. Math. Anal. Appl. 460 (2018), 321--331;
\href{https://arxiv.org/abs/1604.00763}{arXiv:1604.00763}.

\bibitem{BNO}
C. Bargetz, E. A. Nigsch, and N. Ortner,
\emph{A simpler description of the $\kappa$-topologies on the spaces
$\mathcal D_{L^p}$, $L^p$, $\mathcal M^1$},
Math. Nachr. 293 (2020), 1691--1706;
\href{https://arxiv.org/abs/1711.06577}{arXiv:1711.06577}.
\end{thebibliography}

\end{document}
